% @file paper44_group_theory_en.tex
% @brief Paper 44 (EN): Group Theory — Finite Groups, Sylow Theory,
%        Classification, and Applications
% @author Michael Fuhrmann
% @date 2026-03-12
% @build 135

\documentclass[12pt,a4paper]{amsart}
\usepackage{amsmath,amssymb,amsthm}
\usepackage{mathtools}
\usepackage[utf8]{inputenc}
\usepackage[T1]{fontenc}
\usepackage{hyperref}
\usepackage{enumitem,booktabs,array}
\usepackage{geometry}
\geometry{margin=2.5cm}

%% Theorem environments
\newtheorem{theorem}{Theorem}[section]
\newtheorem{lemma}[theorem]{Lemma}
\newtheorem{corollary}[theorem]{Corollary}
\newtheorem{proposition}[theorem]{Proposition}
\newtheorem{conjecture}[theorem]{Conjecture}
\theoremstyle{definition}
\newtheorem{definition}[theorem]{Definition}
\newtheorem{example}[theorem]{Example}
\newtheorem{remark}[theorem]{Remark}

\title{Group Theory: Finite Groups, Sylow Theory,\\
       Classification Theorems, and Applications}

\author{Michael Fuhrmann}
\address{Specialist Mathematics Project, Build 135}
\email{specialist-maths@localhost}
\date{\today}

\begin{document}

\begin{abstract}
Group theory is the mathematical study of symmetry and its algebraic abstraction.
This paper provides a comprehensive treatment of finite group theory, beginning with
the fundamental definitions and the three classical isomorphism theorems. We examine
cyclic and abelian groups, culminating in the complete classification of finite abelian
groups via the Fundamental Theorem. Group actions, the orbit-stabilizer theorem, and
Burnside's lemma are developed in full, followed by the Sylow theorems with complete
proofs. Solvable and nilpotent groups are treated in connection with composition series
and the Jordan--H\"{o}lder theorem, with applications to the insolubility of the quintic
(Abel--Ruffini). The theory of simple groups is discussed, including the Feit--Thompson
theorem and the Classification of Finite Simple Groups (CFSG). The final section surveys
applications in Galois theory, crystallography, and quantum mechanics.
\end{abstract}

\maketitle
\tableofcontents

%% ============================================================
\section{Introduction}
\label{sec:introduction}
%% ============================================================

The concept of a \emph{group} emerged in the early nineteenth century through the work
of \'{E}variste Galois, who introduced what we now call permutation groups in order to
characterise which polynomial equations are solvable by radicals. Independently,
Arthur Cayley defined abstract groups in 1854, recognising that the same algebraic
structure underlies permutations, integers under addition, symmetries of geometric
objects, and a vast range of other mathematical phenomena.

Today group theory is a cornerstone of pure mathematics and theoretical physics.
Its applications span:

\begin{itemize}[leftmargin=2em]
  \item \textbf{Symmetry analysis}: every physical symmetry corresponds to a group;
        conservation laws arise via Noether's theorem from symmetry groups of the action.
  \item \textbf{Galois theory}: the solvability of a polynomial equation by radicals is
        equivalent to the solvability of its Galois group.
  \item \textbf{Cryptography}: the discrete logarithm problem in finite cyclic groups and
        elliptic curve groups underpins modern public-key cryptography.
  \item \textbf{Crystallography}: the 17 wallpaper groups and 230 space groups classify
        all possible crystal symmetries in the plane and in three-dimensional space.
  \item \textbf{Quantum mechanics}: the representation theory of Lie groups such as
        $\mathrm{SU}(2)$ and $\mathrm{SO}(3)$ determines the structure of atomic spectra
        and the spin of elementary particles.
\end{itemize}

This paper is structured as follows. Section~\ref{sec:definitions} lays down the
fundamental definitions. Section~\ref{sec:isomorphism} presents the three isomorphism
theorems. Sections~\ref{sec:cyclic} and~\ref{sec:actions} treat cyclic/abelian groups
and group actions with Sylow theory, respectively. Sections~\ref{sec:solvable}
and~\ref{sec:simple} address solvable/nilpotent groups and simple groups.
Section~\ref{sec:applications} surveys the most important applications.
Section~\ref{sec:conclusion} concludes.

%% ============================================================
\section{Fundamental Definitions}
\label{sec:definitions}
%% ============================================================

\begin{definition}[Group]
A \emph{group} is a pair $(G, \cdot)$ consisting of a non-empty set $G$ together with
a binary operation $\cdot : G \times G \to G$ satisfying:
\begin{enumerate}[label=(\roman*)]
  \item \textbf{Associativity}: $(a \cdot b) \cdot c = a \cdot (b \cdot c)$ for all
        $a,b,c \in G$.
  \item \textbf{Identity}: There exists $e \in G$ such that $e \cdot a = a \cdot e = a$
        for all $a \in G$.
  \item \textbf{Inverses}: For every $a \in G$ there exists $a^{-1} \in G$ such that
        $a \cdot a^{-1} = a^{-1} \cdot a = e$.
\end{enumerate}
If additionally $a \cdot b = b \cdot a$ for all $a, b \in G$, the group is called
\emph{abelian} (or \emph{commutative}). The \emph{order} $|G|$ is the cardinality of $G$.
\end{definition}

\begin{example}
The integers $(\mathbb{Z}, +)$ form an infinite abelian group; the unit circle
$\mathrm{U}(1) = \{z \in \mathbb{C} : |z| = 1\}$ under multiplication is an infinite
abelian group; the symmetric group $S_n$ of all permutations of $\{1, \ldots, n\}$
under composition is a finite group of order $n!$.
\end{example}

\begin{definition}[Subgroup]
A non-empty subset $H \subseteq G$ is a \emph{subgroup} of $G$, written $H \leq G$,
if $H$ is itself a group under the restriction of the operation of $G$. Equivalently,
$H \leq G$ if and only if for all $a, b \in H$ we have $ab^{-1} \in H$.
\end{definition}

\begin{definition}[Coset and Index]
For $H \leq G$ and $a \in G$, the \emph{left coset} of $H$ with respect to $a$ is
$aH = \{ah : h \in H\}$. The number of distinct left cosets is the \emph{index}
$[G : H]$.
\end{definition}

\begin{theorem}[Lagrange, 1771]
\label{thm:lagrange}
Let $G$ be a finite group and $H \leq G$. Then $|H|$ divides $|G|$, and
\[
  |G| = |H| \cdot [G : H].
\]
\end{theorem}

\begin{proof}
The left cosets of $H$ partition $G$ into $[G:H]$ disjoint sets each of size $|H|$.
Counting gives $|G| = [G:H] \cdot |H|$.
\end{proof}

\begin{corollary}
The order of any element $a \in G$ divides $|G|$.
\end{corollary}

\begin{definition}[Normal Subgroup and Quotient Group]
A subgroup $N \leq G$ is \emph{normal}, written $N \trianglelefteq G$, if
$gNg^{-1} = N$ for all $g \in G$. If $N \trianglelefteq G$, the set of cosets
$G/N = \{gN : g \in G\}$ forms a group under $(aN)(bN) = (ab)N$, called the
\emph{quotient group} (or \emph{factor group}) of $G$ by $N$.
\end{definition}

\begin{definition}[Homomorphism and Isomorphism]
A map $\varphi : G \to H$ between groups is a \emph{homomorphism} if
$\varphi(ab) = \varphi(a)\varphi(b)$ for all $a,b \in G$. The \emph{kernel}
of $\varphi$ is $\ker\varphi = \{g \in G : \varphi(g) = e_H\}$ and the
\emph{image} is $\mathrm{im}\,\varphi = \{\varphi(g) : g \in G\}$.
A bijective homomorphism is an \emph{isomorphism}; we write $G \cong H$ if
an isomorphism between $G$ and $H$ exists.
\end{definition}

\begin{proposition}
If $\varphi : G \to H$ is a homomorphism, then $\ker\varphi \trianglelefteq G$
and $\mathrm{im}\,\varphi \leq H$.
\end{proposition}

%% ============================================================
\section{The Isomorphism Theorems}
\label{sec:isomorphism}
%% ============================================================

The three isomorphism theorems are among the most useful structural results in group
theory. They were first stated in their modern form by Emmy Noether in 1927.

\begin{theorem}[First Isomorphism Theorem]
\label{thm:first_iso}
Let $\varphi : G \to H$ be a group homomorphism. Then $\ker\varphi \trianglelefteq G$
and
\[
  G / \ker\varphi \;\cong\; \mathrm{im}\,\varphi.
\]
The isomorphism is given explicitly by $g\ker\varphi \mapsto \varphi(g)$.
\end{theorem}

\begin{proof}
We already know $K := \ker\varphi \trianglelefteq G$.
Define $\tilde\varphi : G/K \to H$ by $\tilde\varphi(gK) = \varphi(g)$.
This is well-defined: if $gK = g'K$ then $g'^{-1}g \in K$, so
$\varphi(g'^{-1}g) = e_H$, giving $\varphi(g) = \varphi(g')$.
It is a homomorphism since $\tilde\varphi((gK)(g'K)) = \varphi(gg') = \varphi(g)\varphi(g')$.
It is injective: $\tilde\varphi(gK) = e_H$ implies $g \in K$, i.e.\ $gK = K = e_{G/K}$.
It is surjective onto $\mathrm{im}\,\varphi$ by construction.
\end{proof}

\begin{remark}
The First Isomorphism Theorem is also called the \emph{Homomorphism Theorem}
(German: \emph{Homomorphiesatz}). It reduces the classification of homomorphic images of
$G$ to the classification of normal subgroups of $G$.
\end{remark}

\begin{theorem}[Second Isomorphism Theorem]
\label{thm:second_iso}
Let $H \leq G$ and $N \trianglelefteq G$. Then $HN = \{hn : h \in H,\, n \in N\}$
is a subgroup of $G$, $H \cap N \trianglelefteq H$, and
\[
  H / (H \cap N) \;\cong\; HN / N.
\]
\end{theorem}

\begin{proof}
One verifies that $HN$ is a subgroup using normality of $N$. The natural map
$\varphi : H \to HN/N$ defined by $\varphi(h) = hN$ is a surjective homomorphism
with kernel $H \cap N$. The result follows from Theorem~\ref{thm:first_iso}.
\end{proof}

\begin{theorem}[Third Isomorphism Theorem]
\label{thm:third_iso}
Let $N \trianglelefteq G$ and $K \trianglelefteq G$ with $K \leq N$. Then
$N/K \trianglelefteq G/K$ and
\[
  (G/K) / (N/K) \;\cong\; G/N.
\]
\end{theorem}

\begin{proof}
The map $\varphi : G/K \to G/N$ defined by $\varphi(gK) = gN$ is a well-defined
surjective homomorphism with kernel $N/K$. Apply Theorem~\ref{thm:first_iso}.
\end{proof}

\begin{corollary}[Correspondence Theorem]
Let $N \trianglelefteq G$. There is a bijective, inclusion-preserving correspondence
between subgroups of $G/N$ and subgroups of $G$ containing $N$. This correspondence
preserves normality.
\end{corollary}

%% ============================================================
\section{Cyclic Groups and Abelian Groups}
\label{sec:cyclic}
%% ============================================================

\begin{definition}[Cyclic Group]
A group $G$ is \emph{cyclic} if there exists $g \in G$ such that every element of
$G$ is a power of $g$: $G = \langle g \rangle = \{g^k : k \in \mathbb{Z}\}$.
\end{definition}

\begin{theorem}
Every cyclic group is abelian. Moreover:
\begin{enumerate}[label=(\roman*)]
  \item Every infinite cyclic group is isomorphic to $(\mathbb{Z}, +)$.
  \item Every finite cyclic group of order $n$ is isomorphic to $\mathbb{Z}/n\mathbb{Z}$.
\end{enumerate}
\end{theorem}

\begin{theorem}
Every subgroup of a cyclic group is cyclic. If $G = \langle g \rangle$ has order $n$,
then the subgroups of $G$ are exactly $\langle g^d \rangle$ for $d \mid n$, and each
subgroup has order $n/d$. In particular, $G$ has exactly one subgroup of each order
dividing $n$.
\end{theorem}

\begin{theorem}[Fundamental Theorem of Finite Abelian Groups]
\label{thm:ftfag}
Every finite abelian group $G$ is isomorphic to a direct product of cyclic groups of
prime-power order:
\[
  G \;\cong\; \mathbb{Z}_{p_1^{a_1}} \times \mathbb{Z}_{p_2^{a_2}} \times \cdots
              \times \mathbb{Z}_{p_k^{a_k}},
\]
where $p_1^{a_1} p_2^{a_2} \cdots p_k^{a_k} = |G|$.
This decomposition is unique up to the ordering of the factors (invariant factor
decomposition).
\end{theorem}

\begin{example}
Up to isomorphism, the abelian groups of order $72 = 2^3 \cdot 3^2$ are:
\begin{align*}
  &\mathbb{Z}_8 \times \mathbb{Z}_9,\quad
   \mathbb{Z}_4 \times \mathbb{Z}_2 \times \mathbb{Z}_9,\quad
   \mathbb{Z}_2^3 \times \mathbb{Z}_9,\\
  &\mathbb{Z}_8 \times \mathbb{Z}_3 \times \mathbb{Z}_3,\quad
   \mathbb{Z}_4 \times \mathbb{Z}_2 \times \mathbb{Z}_3 \times \mathbb{Z}_3,\quad
   \mathbb{Z}_2^3 \times \mathbb{Z}_3^2.
\end{align*}
There are exactly 6 non-isomorphic abelian groups of order 72.
\end{example}

\begin{theorem}[Chinese Remainder Theorem for Groups]
If $\gcd(m,n) = 1$, then $\mathbb{Z}_{mn} \cong \mathbb{Z}_m \times \mathbb{Z}_n$.
\end{theorem}

%% ============================================================
\section{Group Actions and Sylow Theory}
\label{sec:actions}
%% ============================================================

\subsection{Group Actions}

\begin{definition}[Group Action]
A \emph{(left) action} of a group $G$ on a set $X$ is a map
$G \times X \to X$, $(g,x) \mapsto g \cdot x$, satisfying:
\begin{enumerate}[label=(\roman*)]
  \item $e \cdot x = x$ for all $x \in X$,
  \item $(gh) \cdot x = g \cdot (h \cdot x)$ for all $g,h \in G$, $x \in X$.
\end{enumerate}
\end{definition}

\begin{definition}[Orbit and Stabilizer]
For $x \in X$, the \emph{orbit} of $x$ is $G \cdot x = \{g \cdot x : g \in G\}$,
and the \emph{stabilizer} (or isotropy group) of $x$ is
$G_x = \{g \in G : g \cdot x = x\} \leq G$.
\end{definition}

\begin{theorem}[Orbit-Stabilizer Theorem]
\label{thm:orbit_stab}
For any group action of a finite group $G$ on $X$ and any $x \in X$,
\[
  |G \cdot x| \cdot |G_x| = |G|.
\]
\end{theorem}

\begin{theorem}[Burnside's Lemma, 1897]
\label{thm:burnside}
Let $G$ act on a finite set $X$. The number of distinct orbits is
\[
  |X/G| = \frac{1}{|G|} \sum_{g \in G} |X^g|,
\]
where $X^g = \{x \in X : g \cdot x = x\}$ denotes the set of fixed points of $g$.
\end{theorem}

\begin{example}[Necklace Counting]
We count the number of distinct necklaces with 6 beads each coloured one of $c$
colours, where rotations are considered equivalent. The rotation group $G = \mathbb{Z}_6$
acts on the set $X$ of all coloured sequences ($|X| = c^6$).
By Burnside's lemma the number of distinct necklaces is
\[
  N(c) = \frac{1}{6}\bigl(c^6 + c^3 + 2c^2 + 2c\bigr).
\]
For $c = 2$ this gives $N(2) = \tfrac{1}{6}(64 + 8 + 8 + 4) = 14$.
\end{example}

\subsection{Sylow Theory}

The Sylow theorems, proved by Ludwig Sylow in 1872, are the primary tool for understanding
the subgroup structure of finite groups.

\begin{definition}[Sylow $p$-Subgroup]
Let $G$ be a finite group and $p$ a prime with $p^k \mid |G|$ but $p^{k+1} \nmid |G|$.
A subgroup $P \leq G$ with $|P| = p^k$ is called a \emph{Sylow $p$-subgroup} of $G$.
We write $n_p(G)$ for the number of Sylow $p$-subgroups.
\end{definition}

\begin{theorem}[Sylow's First Theorem]
\label{thm:sylow1}
Let $G$ be a finite group, $p$ a prime, and $p^k \mid |G|$. Then $G$ contains a
subgroup of order $p^k$. In particular, Sylow $p$-subgroups exist.
\end{theorem}

\begin{proof}[Proof sketch]
By induction on $|G|$. If $p \nmid [G:Z(G)]$, then $p \mid |Z(G)|$ and Cauchy's theorem
gives an element of order $p$ in the centre. Proceed inductively. In the general case,
use the class equation $|G| = |Z(G)| + \sum [G:C_G(x)]$ and apply induction to a
proper subgroup.
\end{proof}

\begin{theorem}[Sylow's Second Theorem]
\label{thm:sylow2}
All Sylow $p$-subgroups of $G$ are conjugate: for any two Sylow $p$-subgroups $P$
and $Q$ there exists $g \in G$ with $Q = gPg^{-1}$.
\end{theorem}

\begin{theorem}[Sylow's Third Theorem]
\label{thm:sylow3}
The number $n_p = n_p(G)$ of Sylow $p$-subgroups satisfies:
\begin{enumerate}[label=(\roman*)]
  \item $n_p \equiv 1 \pmod{p}$,
  \item $n_p \mid [G : P] = |G|/p^k$.
\end{enumerate}
In particular $n_p = 1$ if and only if the unique Sylow $p$-subgroup is normal in $G$.
\end{theorem}

\begin{example}[Groups of Order 15]
Let $|G| = 15 = 3 \cdot 5$. Then $n_3 \mid 5$ and $n_3 \equiv 1 \pmod{3}$, so
$n_3 = 1$. Similarly $n_5 = 1$. Both Sylow subgroups are normal, so
$G \cong \mathbb{Z}_3 \times \mathbb{Z}_5 \cong \mathbb{Z}_{15}$. There is only one
group of order 15, and it is cyclic.
\end{example}

\begin{corollary}[Cauchy's Theorem]
If $p$ is a prime and $p \mid |G|$, then $G$ contains an element of order $p$.
\end{corollary}

%% ============================================================
\section{Solvable and Nilpotent Groups}
\label{sec:solvable}
%% ============================================================

\subsection{Composition Series and Jordan--H\"{o}lder}

\begin{definition}[Composition Series]
A \emph{composition series} for $G$ is a chain of subgroups
\[
  1 = G_0 \trianglelefteq G_1 \trianglelefteq \cdots \trianglelefteq G_n = G
\]
such that each quotient $G_{i+1}/G_i$ is simple (non-trivial with no proper normal
subgroups). The quotients $G_{i+1}/G_i$ are the \emph{composition factors}.
\end{definition}

\begin{theorem}[Jordan--H\"{o}lder Theorem]
\label{thm:jordan_holder}
If $G$ has a composition series, then any two composition series of $G$ have the same
length, and their composition factors are the same up to permutation and isomorphism.
\end{theorem}

\subsection{Solvable Groups}

\begin{definition}[Derived Series and Solvability]
The \emph{derived series} (or commutator series) of $G$ is
\[
  G = G^{(0)} \geq G^{(1)} \geq G^{(2)} \geq \cdots,
  \quad G^{(i+1)} = [G^{(i)}, G^{(i)}],
\]
where $[H, K]$ denotes the subgroup generated by all commutators $[h,k] = hkh^{-1}k^{-1}$.
The group $G$ is \emph{solvable} (or \emph{soluble}) if $G^{(n)} = 1$ for some $n$.
\end{definition}

\begin{theorem}
A finite group $G$ is solvable if and only if all composition factors in any composition
series of $G$ are cyclic of prime order.
\end{theorem}

\begin{theorem}[Abel--Ruffini Theorem, 1824]
\label{thm:abel_ruffini}
For $n \geq 5$, the symmetric group $S_n$ is \emph{not} solvable (since
$A_n$ is simple and non-abelian). Consequently, a generic polynomial equation of
degree $\geq 5$ cannot be solved by radicals.
\end{theorem}

\begin{proof}[Connection to Galois theory]
A polynomial $f \in \mathbb{Q}[x]$ is solvable by radicals if and only if its Galois
group $\mathrm{Gal}(f)$ is solvable. For the general polynomial of degree $n \geq 5$,
$\mathrm{Gal}(f) \cong S_n$, which contains the simple non-abelian subgroup $A_n$.
Hence $S_n$ is not solvable for $n \geq 5$.
\end{proof}

\subsection{Nilpotent Groups}

\begin{definition}[Lower Central Series and Nilpotency]
The \emph{lower central series} of $G$ is
\[
  G = \gamma_1(G) \geq \gamma_2(G) \geq \cdots,
  \quad \gamma_{i+1}(G) = [G, \gamma_i(G)].
\]
$G$ is \emph{nilpotent} of \emph{class} $c$ if $\gamma_{c+1}(G) = 1$ but
$\gamma_c(G) \neq 1$.
\end{definition}

\begin{theorem}
Every nilpotent group is solvable. A finite group is nilpotent if and only if it is
the direct product of its Sylow $p$-subgroups. In particular, every finite $p$-group
(group of prime power order) is nilpotent.
\end{theorem}

%% ============================================================
\section{Simple Groups}
\label{sec:simple}
%% ============================================================

\begin{definition}[Simple Group]
A group $G$ with $|G| > 1$ is \emph{simple} if its only normal subgroups are $1$ and
$G$ itself.
\end{definition}

\begin{theorem}[$A_n$ is Simple for $n \geq 5$]
\label{thm:an_simple}
The alternating group $A_n$ is simple for all $n \geq 5$.
\end{theorem}

\begin{proof}[Proof sketch]
Every non-trivial normal subgroup of $A_n$ must contain a $3$-cycle (using the fact
that $A_n$ is generated by $3$-cycles and conjugacy class arguments). Any normal
subgroup containing one $3$-cycle contains all $3$-cycles, hence equals $A_n$.
\end{proof}

\subsection{The Projective Special Linear Groups}

\begin{definition}
For a prime power $q = p^r$ and $n \geq 1$, the \emph{projective special linear group}
$\mathrm{PSL}(n, q)$ is the quotient $\mathrm{SL}(n, q) / Z(\mathrm{SL}(n,q))$, where
$\mathrm{SL}(n,q)$ is the group of $n \times n$ matrices over $\mathbb{F}_q$ with
determinant~1.
\end{definition}

\begin{theorem}
For $n \geq 2$ and $q$ a prime power, $\mathrm{PSL}(n, q)$ is simple, except for
$\mathrm{PSL}(2, 2) \cong S_3$ and $\mathrm{PSL}(2, 3) \cong A_4$.
The groups $\mathrm{PSL}(2, q)$ for prime powers $q \geq 4$ are infinite families of
non-abelian simple groups.
\end{theorem}

\subsection{The Feit--Thompson Theorem and CFSG}

\begin{theorem}[Feit--Thompson Theorem, 1963]
\label{thm:feit_thompson}
Every finite group of odd order is solvable.
\end{theorem}

\begin{remark}
The Feit--Thompson theorem occupies an entire volume of the \emph{Pacific Journal of
Mathematics} (255 pages). Its proof is a tour de force of representation theory,
character theory, and combinatorial group theory. As an immediate consequence, every
non-abelian finite simple group has even order.
\end{remark}

\begin{theorem}[Classification of Finite Simple Groups (CFSG), 2004]
\label{thm:cfsg}
Every finite simple group belongs to one of the following families:
\begin{enumerate}[label=(\roman*)]
  \item \textbf{Cyclic groups} $\mathbb{Z}_p$ for prime $p$.
  \item \textbf{Alternating groups} $A_n$ for $n \geq 5$.
  \item \textbf{Groups of Lie type}: 16 infinite families, including
        $\mathrm{PSL}(n,q)$, $\mathrm{PSp}(2n,q)$, $\mathrm{P}\Omega^{\pm}(n,q)$,
        $G_2(q)$, $F_4(q)$, $E_6(q)$, $E_7(q)$, $E_8(q)$, and the twisted variants
        ${}^2A_n(q)$, ${}^2D_n(q)$, ${}^3D_4(q)$, ${}^2G_2(q)$, ${}^2F_4(q)$,
        ${}^2E_6(q)$, ${}^2B_2(q)$.
  \item \textbf{26 sporadic groups}: including the Monster group $\mathbb{M}$
        (order $\approx 8 \times 10^{53}$), the Baby Monster, and the 20
        ``Happy Family'' members related to the Monster via moonshine.
\end{enumerate}
\end{theorem}

\begin{remark}
The proof of the CFSG was completed around 1983 with major contributions by Gorenstein,
Lyons, Solomon, Aschbacher, Smith, and many others. A revised and complete proof,
the ``second generation'' classification, was finished in 2004 in a 12-volume series.
The proof spans over 10,000 pages across hundreds of journal articles.
\end{remark}

\begin{theorem}[Monster Group]
The \emph{Monster group} $\mathbb{M}$ is the largest sporadic simple group. Its order is
\begin{align*}
  |\mathbb{M}| &= 2^{46} \cdot 3^{20} \cdot 5^9 \cdot 7^6 \cdot 11^2 \cdot 13^3
                  \cdot 17 \cdot 19 \cdot 23 \cdot 29 \cdot 31 \cdot 41 \cdot 47
                  \cdot 59 \cdot 71\\
               &\approx 8.08 \times 10^{53}.
\end{align*}
\end{theorem}

%% ============================================================
\section{Applications}
\label{sec:applications}
%% ============================================================

\subsection{Galois Theory}

Group theory and field theory are linked through Galois theory, one of the most beautiful
chapters of mathematics.

\begin{theorem}[Fundamental Theorem of Galois Theory]
\label{thm:galois_fund}
Let $L/K$ be a finite Galois extension with Galois group $\mathrm{Gal}(L/K)$. There
is an inclusion-reversing bijection between intermediate fields $K \subseteq F \subseteq L$
and subgroups $H \leq \mathrm{Gal}(L/K)$:
\[
  F \;\longleftrightarrow\; H = \mathrm{Gal}(L/F), \quad
  [L : F] = |H|, \quad [F : K] = [\mathrm{Gal}(L/K) : H].
\]
Moreover, $F/K$ is Galois if and only if $H \trianglelefteq \mathrm{Gal}(L/K)$, in
which case $\mathrm{Gal}(F/K) \cong \mathrm{Gal}(L/K)/H$.
\end{theorem}

\begin{example}
The splitting field of $x^4 - 2$ over $\mathbb{Q}$ is $L = \mathbb{Q}(\sqrt[4]{2},i)$,
with $\mathrm{Gal}(L/\mathbb{Q}) \cong D_4$ (dihedral group of order 8). The subgroup
lattice of $D_4$ corresponds bijectively to the 10 intermediate fields between
$\mathbb{Q}$ and $L$.
\end{example}

\subsection{Crystallographic Groups}

\begin{definition}
A \emph{space group} is a discrete subgroup $\Gamma$ of the Euclidean group
$\mathrm{E}(3)$ such that the quotient $\mathbb{R}^3/\Gamma$ is compact. The
\emph{point group} is the image of $\Gamma$ in $\mathrm{O}(3)$.
\end{definition}

\begin{theorem}[Crystallographic Restriction Theorem]
In dimension 2, a crystallographic group (wallpaper group) can contain only rotational
symmetries of order $2, 3, 4$, or $6$. There are exactly $17$ wallpaper groups and
exactly $230$ space groups in three dimensions.
\end{theorem}

\begin{remark}
The 230 space groups, classified independently by Fedorov (1891) and Sch\"{o}nflies
(1891), form the mathematical foundation of crystallography. Modern crystal structure
determination via X-ray diffraction exploits this classification at every step.
\end{remark}

\subsection{Quantum Mechanics: \texorpdfstring{$\mathrm{SU}(2)$}{SU(2)} and
            \texorpdfstring{$\mathrm{SO}(3)$}{SO(3)}}

\begin{definition}
The \emph{special unitary group}
$\mathrm{SU}(2) = \{A \in \mathrm{M}_{2}(\mathbb{C}) : A^\dagger A = I,\, \det A = 1\}$
and the \emph{rotation group}
$\mathrm{SO}(3) = \{R \in \mathrm{M}_3(\mathbb{R}) : R^T R = I,\, \det R = 1\}$
are Lie groups related by the double cover $\mathrm{SU}(2) \twoheadrightarrow \mathrm{SO}(3)$
with kernel $\{\pm I\}$.
\end{definition}

\begin{theorem}
The irreducible complex representations of $\mathrm{SU}(2)$ are parametrised by
non-negative half-integers $j \in \{0, \frac{1}{2}, 1, \frac{3}{2}, 2, \ldots\}$
and have dimension $2j+1$. The representations descend to $\mathrm{SO}(3)$ if and only
if $j$ is an integer.
\end{theorem}

\begin{remark}
In quantum mechanics, $j$ is the \emph{spin} quantum number. Electrons, protons, and
neutrons are spin-$\frac{1}{2}$ particles, transforming under the $2$-dimensional
representation of $\mathrm{SU}(2)$. Photons are spin-$1$ particles, transforming under
$\mathrm{SO}(3)$. The Wigner $3j$-symbols that appear in angular momentum coupling are
the Clebsch--Gordan coefficients for the decomposition
$j_1 \otimes j_2 = |j_1 - j_2| \oplus \cdots \oplus (j_1 + j_2)$.
\end{remark}

\subsection{Cryptography: Discrete Logarithm Groups}

\begin{definition}
Let $G = \langle g \rangle$ be a finite cyclic group of order $n$. The
\emph{discrete logarithm problem (DLP)} is: given $h \in G$, find $x \in \mathbb{Z}$
with $g^x = h$.
\end{definition}

\begin{remark}
For carefully chosen groups (e.g.\ prime-order subgroups of $(\mathbb{Z}/p\mathbb{Z})^*$
or points on elliptic curves over finite fields), no polynomial-time classical algorithm
for the DLP is known. This hardness assumption underlies the security of protocols such
as Diffie--Hellman key exchange, ElGamal encryption, and ECDSA signatures. Shor's quantum
algorithm solves the DLP in polynomial time on a quantum computer, motivating the
transition to post-quantum cryptography.
\end{remark}

%% ============================================================
\section{Conclusion}
\label{sec:conclusion}
%% ============================================================

We have surveyed the core of finite group theory: from the basic axioms through the
isomorphism theorems and the classification of finite abelian groups, to the powerful
tools of Sylow theory and the structural results on solvable and nilpotent groups. The
Jordan--H\"{o}lder theorem asserts that composition factors are a group-theoretic
invariant, and the Abel--Ruffini theorem shows how solvability of groups translates
directly into solvability of equations by radicals.

The Classification of Finite Simple Groups (CFSG), one of the greatest collective
achievements in the history of mathematics, provides a complete list of the
``atoms'' from which all finite groups are built. The Feit--Thompson theorem, that
every odd-order finite group is solvable, is a deep prerequisite for the CFSG.

Applications of group theory permeate mathematics and physics. Galois theory uses group
theory to answer ancient questions about polynomial equations. Crystallographic groups
govern the structure of all known crystals. The representation theory of $\mathrm{SU}(2)$
and $\mathrm{SO}(3)$ underpins the quantum theory of angular momentum and spin. Cyclic
groups and their analogues form the backbone of modern public-key cryptography.

Open directions include the moonshine conjectures connecting the Monster group to modular
forms (Borcherds' proof of the Conway--Norton moonshine conjecture, 1992), the Baum--
Connes conjecture relating group $K$-theory to analysis, and the Langlands programme,
which unifies group-theoretic, number-theoretic, and geometric structures through the
representation theory of reductive groups.

%% ============================================================
\section*{References}
%% ============================================================

\begin{thebibliography}{99}

\bibitem{Artin2011}
M.~Artin, \emph{Algebra}, 2nd ed., Pearson, 2011.

\bibitem{Dummit2004}
D.~S.~Dummit and R.~M.~Foote, \emph{Abstract Algebra}, 3rd ed.,
Wiley, Hoboken, NJ, 2004.

\bibitem{Rotman2012}
J.~J.~Rotman, \emph{Advanced Modern Algebra}, 3rd ed., American Mathematical
Society, Providence, RI, 2010.

\bibitem{FeitThompson1963}
W.~Feit and J.~G.~Thompson, ``Solvability of groups of odd order,''
\emph{Pacific J.\ Math.}\ \textbf{13} (1963), no.~3, 775--1029.

\bibitem{Sylow1872}
L.~Sylow, ``Th\'{e}or\`{e}mes sur les groupes de substitutions,''
\emph{Math.\ Ann.}\ \textbf{5} (1872), 584--594.

\bibitem{Gorenstein1982}
D.~Gorenstein, \emph{Finite Simple Groups: An Introduction to Their
Classification}, Plenum Press, New York, 1982.

\bibitem{Conway1985}
J.~H.~Conway and N.~J.~A.~Sloane, \emph{Sphere Packings, Lattices and Groups},
Springer, New York, 1988.

\bibitem{Hall1959}
M.~Hall, \emph{The Theory of Groups}, Macmillan, New York, 1959.

\bibitem{Serre1977}
J.-P.~Serre, \emph{Linear Representations of Finite Groups}, Springer, New York,
1977.

\bibitem{Aschbacher2004}
M.~Aschbacher, R.~Lyons, S.~D.~Smith, and R.~Solomon,
\emph{The Classification of Finite Simple Groups: Groups of Characteristic 2 Type},
Amer.\ Math.\ Soc., 2011.

\end{thebibliography}

\end{document}
